S'observa que els fotons d'una determinada freqüència que incideixen sobre una superfície metàl·lica provoquen l'emissió d'electrons d'aquesta superfície.

\begin{apartat}{1,25}
Si mantenim la intensitat de la llum constant i augmentem la freqüència dels fotons, com varia el nombre d'electrons emesos i l'energia d'aquests? Encercleu la resposta correcta dins del requadre següent i, a continuació, justifiqueu les dues respostes.

\begin{center}
\begin{tabular}{l}
\toprule
Nombre d'electrons emesos: Disminueix / Es manté constant / Augmenta \\
\midrule
Energia dels electrons emesos: Disminueix / Es manté constant / Augmenta \\
\bottomrule
\end{tabular}
\end{center}
\end{apartat}

\begin{apartat}{1,25}
Determineu la funció de treball dels electrons emesos sabent que el potencial de frenada és de 0,29~V, quan la longitud d'ona incident és de 550~nm. Expresseu el resultat en eV.
\end{apartat}

\dades{%
  $1\ \mathrm{eV} = 1,602\times 10^{-19}\ \mathrm{J}$.\\
  $c = 3,00\times 10^{8}\ \mathrm{m\,s^{-1}}$.\\
  $h = 6,63\times 10^{-34}\ \mathrm{J\,s}$.}
